%% GENERATED from PAPER_SUBMISSION.md (pandoc) -- hand-fixed conversion; see BUILD_NOTES.md
\section{Experiments}\label{sec:4}

\subsection{Datasets}\label{sec:4.1}

The primary development and comparator studies use AR2023 5-core categories under the
full-catalog LLOO protocol (\S\ref{sec:3.6}). A separate prospective robustness and
efficiency study uses MovieLens 1M \citep{harper2015movielens} under a frozen global-time
split. AR2023 counts are total 5-core interactions; the MovieLens row reports its
filtered primary view and training-observed catalog:

%% GENERATED by _bestrec_run/emit_latex_tables.py -- DO NOT EDIT BY HAND
% table key: table_datasets41
% provenance: md lines 215-223; family=md-only (pandoc-converted; no JSON family); check={"checked": 0, "note": "md-only table; no JSON family"}
\par\medskip\begingroup\scriptsize\setlength{\tabcolsep}{2pt}\renewcommand{\arraystretch}{1.15}
\noindent
\begin{tabularx}{\linewidth}{>{\raggedright\arraybackslash}p{0.18\linewidth}>{\raggedright\arraybackslash}p{0.29\linewidth}>{\raggedleft\arraybackslash}p{0.09\linewidth}>{\raggedleft\arraybackslash}p{0.12\linewidth}>{\raggedleft\arraybackslash}p{0.20\linewidth}}
\toprule
Category & Role in this paper & Users & Items & Interactions (total) \\
\midrule
\textbf{Video\_\allowbreak Games} & descriptive system context (\S\ref{sec:5.1}); tail-pattern null (equivalence claim retracted \S\ref{sec:5.3}) & 94,762 & 25,612 & 814,586 \\
\textbf{Musical\_\allowbreak Instruments} & outcome-known FIR internal study + pre-declared per-category point-estimate comparison (\S\ref{sec:5.2}) & 57,439 & 24,587 & 511,836 \\
\textbf{Office\_\allowbreak Products} & V1 prereg \textbf{VOID} (descriptive; Appendix A.0); \textbf{redesigned V3 prereg PASSED} (\S\ref{sec:5.2}; \texttt{OFFICE\_V3\_RESULTS.md}) & 223,308 & 77,551 & 1,800,878 \\
\textbf{Beauty\_and\_\allowbreak Personal\_Care} & tail-pattern exploratory 3-seed estimate (\S\ref{sec:5.3}); superseded cross-pipeline scan (Supplement S.1) & 729,576 & 207,649 & 6,624,441 \\
Industrial\_and\_\allowbreak Scientific & \textbf{pre-declared FIR-breadth: artifact-PASS} (frozen rule fired; paired premise withdrawn; post-hoc independent-arm estimate --- \S\ref{sec:5.2}; \texttt{FIR\_BREADTH\_RESULTS.md}) & 50,985 & 25,848 & 412,947 \\
CDs\_and\_\allowbreak Vinyl & \textbf{pre-declared FIR-breadth: artifact-PASS} (frozen rule fired; paired premise withdrawn; post-hoc independent-arm estimate --- \S\ref{sec:5.2}; \texttt{FIR\_BREADTH\_RESULTS.md}) & 123,876 & 89,370 & 1,552,764 \\
MovieLens 1M, rating\ensuremath{\ge}4 & \textbf{prospectively frozen non-Amazon FIR replication/efficiency study; negative verdict} (\S\ref{sec:5.2}; Harper \& Konstan, 2015) & 1,033 & 2,359 training-observed & 575,281 filtered events; 135,131 train rows \\
\bottomrule
\end{tabularx}
\endgroup\par\medskip


Video\_Games/Musical\_Instruments/Office\_Products match the HSTU-BLaIR comparator pipeline's own statistics exactly on users/items (interactions within \ensuremath{\pm}1; \S\ref{sec:5.2}). The Beauty\_and\_\allowbreak PC catalog is 8\ensuremath{\times} larger than Video\_Games, with median user history of only 5 interactions --- a harder benchmark. Every pre-declared campaign that has completed is mechanically adjudicated (verdicts at the relevant result blocks: \S\ref{sec:5.2}, \S\ref{sec:5.7}, \S\ref{sec:5.8}, and Appendix~A.0).


\subsection{Baselines}\label{sec:4.2}

\begin{itemize}
\tightlist
\item
  Popularity (trivial floor)
\item
  SASRec with sampled-512 softmax (our ``original'' baseline)
\item
  SASRec with chunked-full-softmax (improved baseline)
\item
  BERT4Rec \citep{sun2019bert4rec}
\item
  Published comparators: TIGER \citep{rajput2023tiger}, BLaIR \citep{hou2024blair}, LIGER \citep{yang2024liger}
\end{itemize}


\subsection{Experimental design (headline runs)}\label{sec:4.3}

All headline results (\S\ref{sec:5.1}--\S\ref{sec:5.4}) use the HSTU-style pure-PyTorch encoder (\S\ref{sec:3.2}) under the fixed AR2023 5-core full-catalog LLOO protocol (\S\ref{sec:3.6}). Each configuration is trained for \textbf{40 epochs} (the \S\ref{sec:5.1} Video\_Games headline; the pre-declared-file MI/Office V3/FIR-breadth campaigns train 20 epochs per their frozen configurations) at batch size 256 (d\_model 64, 4 layers, 2 heads, dropout 0.5) with the memory-efficient chunked-full-softmax loss (\S\ref{sec:3.5}, item chunk 32,768) and a warmup-cosine learning-rate schedule. On top of the bias stack (TAPE-512 + time bias + text-similarity bias + pos-rab) the winning configuration adds label smoothing (\ensuremath{\varepsilon}=0.2) and the causal FIR filter (K=8). We run \textbf{6 seeds (20260608--20260613)} for the full-model headline and \textbf{5 seeds (20260608--20260612)} for every per-component ablation rung, reporting mean \ensuremath{\pm} sample-std; the reported test number is always selected by best validation NDCG@10 (best-by-val), never by test. \textbf{Fixed-split inference scope (stated prominently, 2026-07-20):} the training seed is the only randomized unit; every interval in this paper quantifies optimization variability on one fixed public split per category --- nothing here estimates user-resampling, split-choice, category-sampling, or cross-dataset uncertainty, and all claims are scoped to these exact datasets and splits. Evaluation is full-catalog (n\_eval = 94,762 on Video\_Games), with all train+val items masked (\S\ref{sec:3.6}).

The tail analysis (\S\ref{sec:5.3}--\S\ref{sec:5.4}) rests on two arm comparisons, each toggling one named configuration flag on the same split and seed numbers (bundled interventions: the arms are not initialization-paired and optimizer paths differ) and reporting \texttt{by\_popularity} NDCG@10 / HR@10 on leak-free train-frequency terciles (terciles frozen at full density):

\begin{itemize}
\tightlist
\item
  \textbf{text vs ID-only:} the full text stack against an ID-only ablation (\texttt{-\/-no-sbert}, no prototypes, no text-similarity bias).
\item
  \textbf{density / connectivity titration (\S\ref{sec:5.4}):} training-only sub-sampling that thins either interactions (interaction-mode) or whole users (user-mode) to match a sparser reference category's resource levels, with the evaluation set held fixed.
\end{itemize}

A separate \textbf{20-variant cross-pipeline transfer scan on Beauty\_and\_\allowbreak PC} --- supporting reproducibility material, not part of the headline spine --- is reported in Appendix A.1.

The MovieLens study used six exact matched-initialization arms over eight pre-specified
seed blocks: identity, channel-shared K=16, eight-group K=16, rank-at-most-8
tangent-factorized K=16, per-channel K=16, and an equal-parameter current-position-only
placebo. Acquisition had already read the official ratings file, constructed and
hash-bound the TEST split, and used target-in-training-catalog eligibility to define
the primary cohort. During fitting and validation selection, trainers consumed
TRAIN+VALID only and emitted no TEST scores. After all 96 training runs completed,
each selected checkpoint received one sealed TEST evaluation. The primary estimand was NDCG@10 on rating-at-least-4 events
after a global 90\% time boundary; all ratings, user/item cluster bootstraps, and
resource readings were frozen sensitivities. The three-candidate noninferiority family
used margin 0.000500 and Holm-adjusted one-sided tests.


\subsection{Training details and hardware}\label{sec:4.4}

Headline runs use Adam (lr 1e-3, weight decay 1e-5) with gradient clipping (norm 5.0) under the 40-epoch warmup-cosine schedule above. Hardware: a single NVIDIA RTX 5060 Ti (16 GB, Blackwell sm\_120). Video\_Games training takes \textasciitilde10 min/seed; Beauty\_and\_\allowbreak PC \textasciitilde1.5--2 hr/seed with the chunked-full-softmax loss. (The Appendix A.1 Beauty scan used a shorter 15--20-epoch budget and seeds 42/43; those details are local to that supporting study.)

